% --- Archon physics formalization source begin ---
% archon:physics
% archon:covers PhyXMiniProblems/problem_phyx_mini_0681.lean
% archon:source-report reports/phyx_mini/problem_phyx_mini_0681.source.json

\chapter{Physics problem phyx\_mini\_0681}
\label{ch:PhyXMiniProblems_problem_phyx_mini_0681}

\paragraph{Problem source.}
\#\# Physical scenario

In figure a spider is resting after starting to spin its web. The gravitational force on the spider makes it exert a downward force of \textbackslash{}( 0.150 \textbackslash{}text\{ N\} \textbackslash{}) on the junction of the three strands of silk. The junction is supported by different tension forces in the two strands above it so that the resultant force on the junction is zero. The two sloping strands are perpendicular, and we have chosen the \textbackslash{}( x \textbackslash{}) and \textbackslash{}( y \textbackslash{}) directions to be along them. The tension \textbackslash{}( T\_x \textbackslash{}) is \textbackslash{}( 0.127 \textbackslash{}text\{ N\} \textbackslash{}).

\#\# Question

Find the tension \textbackslash{}( T\_y \textbackslash{}

\#\# Answer choices

- A: A: \textbackslash{}( 0.039 \textbackslash{}text\{ N\} \textbackslash{})

- B: B: \textbackslash{}( 0.029 \textbackslash{}text\{ N\} \textbackslash{})

- C: C: \textbackslash{}( 0.078 \textbackslash{}text\{ N\} \textbackslash{})

- D: D: \textbackslash{}( 0.156 \textbackslash{}text\{ N\} \textbackslash{})

\#\# Figure caption (auxiliary; use the image as primary evidence)

The image shows a spider hanging from a web, positioned at the corner of a ceiling. The corner is illustrated using wooden beams. At the intersecting point of the ceiling, two tension components are defined:

1. **Ty**: A tension force along the vertical (y-axis) direction.
2. **Tx**: A tension force along the horizontal (x-axis) direction.

The axes (x and y) are indicated with arrows, with the origin placed at the corner where the two tension components meet. The spider is hanging directly downward from this point. The illustration emphasizes mechanical forces acting on the spider in a physics context.

\#\# Dataset metadata

Statics; Spatial Relation Reasoning, Multi-Formula Reasoning

\paragraph{Recorded answer/context.}
C: C: \textbackslash{}( 0.078 \textbackslash{}text\{ N\} \textbackslash{})

\paragraph{Figure/image path.}
/root/proposal\_for\_physic/science-mango/phyx\_mini\_run/phyx\_data/test\_image/681.png

\paragraph{Formalization target.}
create a compiling Lean file with sorry bodies at `PhyXMiniProblems/problem\_phyx\_mini\_0681.lean`. The Lean declarations must preserve the physical quantities, dimensions or dimensional roles, figure labels, governing-law hypotheses, and final relation expressed by this problem.
Use Mathlib/Physlib names found through LeanExplore where available. If a physics API is missing, introduce faithful local abstractions rather than scalar placeholder aliases.

\begin{theorem}[Physics formalization target]
\label{thm:physics:phyx_mini_0681:target}
\lean{PhyXMiniProblems.ProblemPhyXMini0681.problem_phyx_mini_0681}
\uses{def:physics:phyx-mini-0681:phyxminiproblems-problemphyxmini0681-forcemagnitudeinnewtons, def:physics:phyx-mini-0681:phyxminiproblems-problemphyxmini0681-spiderwebstaticssetup, def:physics:phyx-mini-0681:phyxminiproblems-problemphyxmini0681-matchessuppliedspiderwebfigure, def:physics:phyx-mini-0681:phyxminiproblems-problemphyxmini0681-matchesproblemscenario, def:physics:phyx-mini-0681:phyxminiproblems-problemphyxmini0681-matchesproblemforcereadouts, def:physics:phyx-mini-0681:phyxminiproblems-problemphyxmini0681-hasperpendicularslopingstrandaxes, def:physics:phyx-mini-0681:phyxminiproblems-problemphyxmini0681-obeysidealsilktensiondirectionlaw, def:physics:phyx-mini-0681:phyxminiproblems-problemphyxmini0681-obeysjunctionforcebalance, def:physics:phyx-mini-0681:phyxminiproblems-problemphyxmini0681-isclosestdisplayedanswer}
The assigned autoformalize agent should translate this physics problem into Lean declarations in the covered file, with theorem and lemma proof bodies written as `by sorry`.
\end{theorem}
\begin{proof}
This is an autoformalization task, not a proof task. Produce statements that can later be proved by the physics prover without weakening the physical model.
\end{proof}

% --- Archon declaration topology begin ---
\section{Declaration topology}
\begin{definition}[force Dimension]
  \label{def:physics:phyx-mini-0681:phyxminiproblems-problemphyxmini0681-forcedimension}
  \lean{PhyXMiniProblems.ProblemPhyXMini0681.forceDimension}
  The physical dimension M L T⁻² of force
\end{definition}
\begin{proof}
  This is the declared model definition of force Dimension.
\end{proof}
\begin{definition}[Planar Force Quantity]
  \label{def:physics:phyx-mini-0681:phyxminiproblems-problemphyxmini0681-planarforcequantity}
  \lean{PhyXMiniProblems.ProblemPhyXMini0681.PlanarForceQuantity}
  \uses{def:physics:phyx-mini-0681:phyxminiproblems-problemphyxmini0681-forcedimension}
  A unit-independent planar physical force quantity
\end{definition}
\begin{proof}
  This is the declared model definition of Planar Force Quantity.
\end{proof}
\begin{definition}[Diagram Vector]
  \label{def:physics:phyx-mini-0681:phyxminiproblems-problemphyxmini0681-diagramvector}
  \lean{PhyXMiniProblems.ProblemPhyXMini0681.DiagramVector}
  A dimensionless vector used for directions and locations in the diagram
\end{definition}
\begin{proof}
  This is the declared model definition of Diagram Vector.
\end{proof}
\begin{definition}[downward Diagram Direction]
  \label{def:physics:phyx-mini-0681:phyxminiproblems-problemphyxmini0681-downwarddiagramdirection}
  \lean{PhyXMiniProblems.ProblemPhyXMini0681.downwardDiagramDirection}
  \uses{def:physics:phyx-mini-0681:phyxminiproblems-problemphyxmini0681-diagramvector}
  Downward in diagram coordinates: coordinate 0 is right and 1 is up
\end{definition}
\begin{proof}
  This is the declared model definition of downward Diagram Direction.
\end{proof}
\begin{definition}[force Vector In Newtons]
  \label{def:physics:phyx-mini-0681:phyxminiproblems-problemphyxmini0681-forcevectorinnewtons}
  \lean{PhyXMiniProblems.ProblemPhyXMini0681.forceVectorInNewtons}
  \uses{def:physics:phyx-mini-0681:phyxminiproblems-problemphyxmini0681-planarforcequantity}
  Read a planar force vector in coherent SI units, i.e. newtons
\end{definition}
\begin{proof}
  This is the declared model definition of force Vector In Newtons.
\end{proof}
\begin{definition}[force Magnitude In Newtons]
  \label{def:physics:phyx-mini-0681:phyxminiproblems-problemphyxmini0681-forcemagnitudeinnewtons}
  \lean{PhyXMiniProblems.ProblemPhyXMini0681.forceMagnitudeInNewtons}
  \uses{def:physics:phyx-mini-0681:phyxminiproblems-problemphyxmini0681-planarforcequantity, def:physics:phyx-mini-0681:phyxminiproblems-problemphyxmini0681-forcevectorinnewtons}
  The Euclidean magnitude of a coherent-SI force readout, in newtons
\end{definition}
\begin{proof}
  This is the declared model definition of force Magnitude In Newtons.
\end{proof}
\begin{definition}[Sloping Axis]
  \label{def:physics:phyx-mini-0681:phyxminiproblems-problemphyxmini0681-slopingaxis}
  \lean{PhyXMiniProblems.ProblemPhyXMini0681.SlopingAxis}
  The two sloping strands, simultaneously used as the labelled coordinate axes
\end{definition}
\begin{proof}
  This is the declared model definition of Sloping Axis.
\end{proof}
\begin{definition}[Figure Direction]
  \label{def:physics:phyx-mini-0681:phyxminiproblems-problemphyxmini0681-figuredirection}
  \lean{PhyXMiniProblems.ProblemPhyXMini0681.FigureDirection}
  Qualitative directions of lines and arrows in the supplied raster
\end{definition}
\begin{proof}
  This is the declared model definition of Figure Direction.
\end{proof}
\begin{definition}[expected Upper Strand Direction]
  \label{def:physics:phyx-mini-0681:phyxminiproblems-problemphyxmini0681-expectedupperstranddirection}
  \lean{PhyXMiniProblems.ProblemPhyXMini0681.expectedUpperStrandDirection}
  \uses{def:physics:phyx-mini-0681:phyxminiproblems-problemphyxmini0681-slopingaxis, def:physics:phyx-mini-0681:phyxminiproblems-problemphyxmini0681-figuredirection}
  The direction in which each upper strand leaves the junction in the figure
\end{definition}
\begin{proof}
  This is the declared model definition of expected Upper Strand Direction.
\end{proof}
\begin{definition}[Figure Feature]
  \label{def:physics:phyx-mini-0681:phyxminiproblems-problemphyxmini0681-figurefeature}
  \lean{PhyXMiniProblems.ProblemPhyXMini0681.FigureFeature}
  Named physical objects and silk segments visible in image 681.png
\end{definition}
\begin{proof}
  This is the declared model definition of Figure Feature.
\end{proof}
\begin{definition}[Figure Label]
  \label{def:physics:phyx-mini-0681:phyxminiproblems-problemphyxmini0681-figurelabel}
  \lean{PhyXMiniProblems.ProblemPhyXMini0681.FigureLabel}
  Literal labels printed in the supplied raster
\end{definition}
\begin{proof}
  This is the declared model definition of Figure Label.
\end{proof}
\begin{definition}[Spider Web Figure]
  \label{def:physics:phyx-mini-0681:phyxminiproblems-problemphyxmini0681-spiderwebfigure}
  \lean{PhyXMiniProblems.ProblemPhyXMini0681.SpiderWebFigure}
  \uses{def:physics:phyx-mini-0681:phyxminiproblems-problemphyxmini0681-slopingaxis, def:physics:phyx-mini-0681:phyxminiproblems-problemphyxmini0681-figuredirection, def:physics:phyx-mini-0681:phyxminiproblems-problemphyxmini0681-figurefeature, def:physics:phyx-mini-0681:phyxminiproblems-problemphyxmini0681-figurelabel}
  Typed qualitative content of the primary figure. The bitmap supplies the strand geometry and the x, y, Tₓ, and Tᵧ labels, but it contains no numerical force readout
\end{definition}
\begin{proof}
  This is the declared model definition of Spider Web Figure.
\end{proof}
\begin{definition}[Spider Web Statics Setup]
  \label{def:physics:phyx-mini-0681:phyxminiproblems-problemphyxmini0681-spiderwebstaticssetup}
  \lean{PhyXMiniProblems.ProblemPhyXMini0681.SpiderWebStaticsSetup}
  \uses{def:physics:phyx-mini-0681:phyxminiproblems-problemphyxmini0681-planarforcequantity, def:physics:phyx-mini-0681:phyxminiproblems-problemphyxmini0681-diagramvector, def:physics:phyx-mini-0681:phyxminiproblems-problemphyxmini0681-slopingaxis, def:physics:phyx-mini-0681:phyxminiproblems-problemphyxmini0681-spiderwebfigure}
  Independent physical quantities and geometry for the web junction. In particular, the y-strand tension and resultant are fields to be constrained by the governing laws, not definitions made from the recorded answer
\end{definition}
\begin{proof}
  This is the declared model definition of Spider Web Statics Setup.
\end{proof}
\begin{definition}[Matches Supplied Spider Web Figure]
  \label{def:physics:phyx-mini-0681:phyxminiproblems-problemphyxmini0681-matchessuppliedspiderwebfigure}
  \lean{PhyXMiniProblems.ProblemPhyXMini0681.MatchesSuppliedSpiderWebFigure}
  \uses{def:physics:phyx-mini-0681:phyxminiproblems-problemphyxmini0681-expectedupperstranddirection, def:physics:phyx-mini-0681:phyxminiproblems-problemphyxmini0681-spiderwebstaticssetup}
  Literal evidence transcribed from the primary raster: two differently sloped upper strands meet the vertical spider strand at the labelled axis origin
\end{definition}
\begin{proof}
  This is the declared model definition of Matches Supplied Spider Web Figure.
\end{proof}
\begin{definition}[Matches Problem Scenario]
  \label{def:physics:phyx-mini-0681:phyxminiproblems-problemphyxmini0681-matchesproblemscenario}
  \lean{PhyXMiniProblems.ProblemPhyXMini0681.MatchesProblemScenario}
  \uses{def:physics:phyx-mini-0681:phyxminiproblems-problemphyxmini0681-downwarddiagramdirection, def:physics:phyx-mini-0681:phyxminiproblems-problemphyxmini0681-forcevectorinnewtons, def:physics:phyx-mini-0681:phyxminiproblems-problemphyxmini0681-forcemagnitudeinnewtons, def:physics:phyx-mini-0681:phyxminiproblems-problemphyxmini0681-spiderwebstaticssetup}
  All three strand forces are concurrent at the junction, and the vertical spider strand transmits a downward force to that junction
\end{definition}
\begin{proof}
  This is the declared model definition of Matches Problem Scenario.
\end{proof}
\begin{definition}[Matches Problem Force Readouts]
  \label{def:physics:phyx-mini-0681:phyxminiproblems-problemphyxmini0681-matchesproblemforcereadouts}
  \lean{PhyXMiniProblems.ProblemPhyXMini0681.MatchesProblemForceReadouts}
  \uses{def:physics:phyx-mini-0681:phyxminiproblems-problemphyxmini0681-forcemagnitudeinnewtons, def:physics:phyx-mini-0681:phyxminiproblems-problemphyxmini0681-spiderwebstaticssetup}
  The two numerical force magnitudes stated in the prose. The downward load is 0.150 N and the supplied x-strand tension is 0.127 N; no Tᵧ value or answer label occurs in these data
\end{definition}
\begin{proof}
  This is the declared model definition of Matches Problem Force Readouts.
\end{proof}
\begin{definition}[Has Perpendicular Sloping Strand Axes]
  \label{def:physics:phyx-mini-0681:phyxminiproblems-problemphyxmini0681-hasperpendicularslopingstrandaxes}
  \lean{PhyXMiniProblems.ProblemPhyXMini0681.HasPerpendicularSlopingStrandAxes}
  \uses{def:physics:phyx-mini-0681:phyxminiproblems-problemphyxmini0681-spiderwebstaticssetup}
  The two diagram axes are unit directions along perpendicular silk strands
\end{definition}
\begin{proof}
  This is the declared model definition of Has Perpendicular Sloping Strand Axes.
\end{proof}
\begin{definition}[Obeys Ideal Silk Tension Direction Law]
  \label{def:physics:phyx-mini-0681:phyxminiproblems-problemphyxmini0681-obeysidealsilktensiondirectionlaw}
  \lean{PhyXMiniProblems.ProblemPhyXMini0681.ObeysIdealSilkTensionDirectionLaw}
  \uses{def:physics:phyx-mini-0681:phyxminiproblems-problemphyxmini0681-forcevectorinnewtons, def:physics:phyx-mini-0681:phyxminiproblems-problemphyxmini0681-forcemagnitudeinnewtons, def:physics:phyx-mini-0681:phyxminiproblems-problemphyxmini0681-spiderwebstaticssetup}
  Ideal silk tension pulls away from the junction along its strand. This is a general tension-direction law and does not prescribe either tension magnitude
\end{definition}
\begin{proof}
  This is the declared model definition of Obeys Ideal Silk Tension Direction Law.
\end{proof}
\begin{definition}[Obeys Junction Force Balance]
  \label{def:physics:phyx-mini-0681:phyxminiproblems-problemphyxmini0681-obeysjunctionforcebalance}
  \lean{PhyXMiniProblems.ProblemPhyXMini0681.ObeysJunctionForceBalance}
  \uses{def:physics:phyx-mini-0681:phyxminiproblems-problemphyxmini0681-forcevectorinnewtons, def:physics:phyx-mini-0681:phyxminiproblems-problemphyxmini0681-spiderwebstaticssetup}
  Vector superposition at the junction and static equilibrium. The spider load is the force transmitted downward through the third strand. Neither field states the requested Tᵧ magnitude or a displayed answer
\end{definition}
\begin{proof}
  This is the declared model definition of Obeys Junction Force Balance.
\end{proof}
\begin{lemma}[perpendicular Tension Magnitude Relation]
  \label{lem:physics:phyx-mini-0681:phyxminiproblems-problemphyxmini0681-perpendiculartensionmagnituderelation}
  \lean{PhyXMiniProblems.ProblemPhyXMini0681.perpendicularTensionMagnitudeRelation}
  \uses{def:physics:phyx-mini-0681:phyxminiproblems-problemphyxmini0681-forcemagnitudeinnewtons, def:physics:phyx-mini-0681:phyxminiproblems-problemphyxmini0681-spiderwebstaticssetup, def:physics:phyx-mini-0681:phyxminiproblems-problemphyxmini0681-hasperpendicularslopingstrandaxes, def:physics:phyx-mini-0681:phyxminiproblems-problemphyxmini0681-obeysidealsilktensiondirectionlaw, def:physics:phyx-mini-0681:phyxminiproblems-problemphyxmini0681-obeysjunctionforcebalance}
  Perpendicular tension directions and zero resultant make the squared spider load equal to the sum of the squared upper-strand tensions
\end{lemma}
\begin{proof}
  The conclusion follows from the governing relations and figure readouts named in the statement of perpendicular Tension Magnitude Relation.
\end{proof}
\begin{definition}[Answer Choice]
  \label{def:physics:phyx-mini-0681:phyxminiproblems-problemphyxmini0681-answerchoice}
  \lean{PhyXMiniProblems.ProblemPhyXMini0681.AnswerChoice}
  The four answer labels printed with the problem
\end{definition}
\begin{proof}
  This is the declared model definition of Answer Choice.
\end{proof}
\begin{definition}[displayed Answer Tension In Newtons]
  \label{def:physics:phyx-mini-0681:phyxminiproblems-problemphyxmini0681-displayedanswertensioninnewtons}
  \lean{PhyXMiniProblems.ProblemPhyXMini0681.displayedAnswerTensionInNewtons}
  \uses{def:physics:phyx-mini-0681:phyxminiproblems-problemphyxmini0681-answerchoice}
  The Tᵧ magnitude in newtons printed beside each answer label
\end{definition}
\begin{proof}
  This is the declared model definition of displayed Answer Tension In Newtons.
\end{proof}
\begin{definition}[Is Closest Displayed Answer]
  \label{def:physics:phyx-mini-0681:phyxminiproblems-problemphyxmini0681-isclosestdisplayedanswer}
  \lean{PhyXMiniProblems.ProblemPhyXMini0681.IsClosestDisplayedAnswer}
  \uses{def:physics:phyx-mini-0681:phyxminiproblems-problemphyxmini0681-answerchoice, def:physics:phyx-mini-0681:phyxminiproblems-problemphyxmini0681-displayedanswertensioninnewtons}
  A choice is closest to an exact tension when no displayed magnitude has a smaller absolute error. This records multiple-choice selection without identifying the exact physical value with a rounded or inconsistent display
\end{definition}
\begin{proof}
  This is the declared model definition of Is Closest Displayed Answer.
\end{proof}
% --- Archon declaration topology end ---
% --- Archon physics formalization source end ---
